\documentclass[11pt,a4paper]{article}

%% =====================================================================
%% Clay Substrate-Reformulation Paper --- HONEST-SCOPE REWRITE
%%
%% Title: A Substrate-Level Reformulation of the Six Remaining Clay
%%        Millennium Problems via the Principia Fractalis Framework
%%        (with alpha_Poincare = 1 as the Substrate's Editorial
%%         Convention Motivated by Perelman 2003)
%%
%% Author: Pablo Cohen (with Claude Opus 4.7 / Anthropic)
%% Date  : 2026-06-06
%% Anchor: HEAD 2b777ca of FractalDevTeam/Principia-Fractalis.
%%         PF_Lean4_Code: kernel-only axioms
%%         [propext, Classical.choice, Quot.sound].
%%         PF_Lean4Lean: same-language Path-C second pass through
%%         Lean's kernel from a separate package; NOT an independent
%%         type-checker.
%% Honest-scope discipline:
%%         Every claim in this paper is sourceable to a specific Lean
%%         file's own HONEST SCOPE section. Where a Lean file flags a
%%         field as :=True or := rfl or substrate-shadow, the paper
%%         flags the same. Earlier drafts drifted into Clay-prize
%%         framings; this rev corrects the drift on the prose side.
%% =====================================================================

\usepackage{amsmath,amsthm,amssymb,amsfonts}
\usepackage[margin=1in]{geometry}
\usepackage{hyperref}
\usepackage{cleveref}
\usepackage{microtype}
\usepackage{enumitem}
\usepackage{booktabs}
\usepackage{xcolor}

\hypersetup{
  colorlinks=true,
  linkcolor=blue!50!black,
  urlcolor=blue!50!black,
  citecolor=blue!50!black,
}

\theoremstyle{plain}
\newtheorem{theorem}{Theorem}[section]
\newtheorem{lemma}[theorem]{Lemma}
\newtheorem{proposition}[theorem]{Proposition}
\newtheorem{corollary}[theorem]{Corollary}

\theoremstyle{definition}
\newtheorem{definition}[theorem]{Definition}
\newtheorem{remark}[theorem]{Remark}

\title{A Substrate-Level Reformulation of the Six Remaining
Clay Millennium Problems via the Principia Fractalis Framework\\
\large (with $\alpha_{\mathrm{Poincar\acute e}} = 1$ as the
Substrate's Editorial Convention Motivated by Perelman 2003)}
\author{Pablo Cohen\\
\small{\texttt{psolorzano@gmail.com}}\\[0.3em]
\small{Collaboration: Claude Opus 4.7 (Anthropic)}}
\date{2026-06-06}

\begin{document}
\maketitle

\begin{abstract}
This paper presents a \emph{substrate-level reformulation} of the
six remaining Clay Millennium Problems (Riemann Hypothesis,
$\mathsf{P}$ vs $\mathsf{NP}$, Navier--Stokes, Yang--Mills,
Birch--Swinnerton-Dyer, Hodge) via the Principia Fractalis
framework. \textbf{This is not a Clay-prize claim.} The six
remaining axes are \emph{not} discharged at the literal mathlib
tier required by the Clay rules. The mechanism we expose is the
following: (i) the framework's $\alpha$-skeleton, with the
editorial convention $\alpha_{\mathrm{Poincar\acute e}} := 1$
motivated by Perelman 2003's resolution of the Poincar\'e
conjecture~\cite{perelman2002,perelman2003a,perelman2003b}; (ii)
eleven \emph{framework-internal arithmetic compatibility
identities} on those nine chosen $\alpha$-values (machine-verified
in $\mathtt{PF.CrossMillenniumSharedInvariants}$ but explicitly
flagged in that file as ``\emph{not Millennium discharges} ---
they are axiom-free algebraic facts about the 9 chosen
$\alpha$-values; the framework's interpretation is conjectural''
\cite[file header]{pf_lean_code}); (iii) canonical encodings of
each Clay axis on the framework's own carriers (which differ from
the literal mathlib carriers in named, disclosed ways); (iv) a
simultaneous-closure theorem that takes the editorial convention
plus a seven-field residual bundle and produces all six
$\mathtt{Clay\_*\_Standard}$ contracts \emph{on the framework's
canonical encodings} (machine-verified at HEAD \texttt{2b777ca},
kernel axioms $\mathtt{[propext, Classical.choice, Quot.sound]}$).
The simultaneous-closure mechanism on the substrate is the
genuine contribution; we submit it for evaluation as a
structurally illuminating reformulation worth editorial attention
as a research direction. Independent expert review and
substantial revision are explicitly invited.
\end{abstract}

\tableofcontents

%% =====================================================================
\section{The framework's claim}
\label{sec:claim}

\subsection{What this paper offers and what it does not offer}
\label{sec:offers}

This paper offers a \emph{substrate-level reformulation} of the
six remaining Clay Millennium Problems, machine-verified in
Lean~4 with kernel-only axioms. It does \emph{not} offer a
literal mathlib-tier Clay-prize discharge for any of them. We
foreground this distinction in the opening section because the
file-level honest-scope discipline of the supporting Lean tree
demands it, and because earlier drafts of this paper drifted
into framings that overstated what the Lean kernel carries.

\medskip

\noindent\fbox{\parbox{0.97\textwidth}{
\textbf{What the Lean kernel carries.} A simultaneous-closure
theorem on the framework's own canonical encodings, gated by
(i) the editorial convention $\alpha_{\mathrm{Poincar\acute e}}
:= 1$, and (ii) a seven-field bundle of which two fields are
$\mathtt{True}$-markers and five are typed residuals. The
discharges land on the framework's substrate-level encodings,
which differ from the literal mathlib carriers in disclosed,
named ways (substrate-shadow vs.\ literal Chow group; finite-dim
$\mathtt{Fin}\,2 \to \mathbb{R}$ vs.\ inf-dim $L^2$; six-piece
framework-Prop conjunction vs.\ literal PDE; et cetera).

\textbf{What this paper does \emph{not} claim.} A discharge of
Riemann, P vs.\ NP, Navier--Stokes, Yang--Mills, BSD, or Hodge
at the literal mathlib tier that the Clay rules require. The
six remaining axes remain open at that tier; the framework's
contribution is structural reformulation, not prize-eligible
discharge.
}}

\subsection{The editorial convention $\alpha_{\mathrm{Poincar\acute e}} = 1$}
\label{sec:editorial-convention}

The framework adopts $\alpha_{\mathrm{Poincar\acute e}} := 1$ as
its editorial convention at the Poincar\'e axis of the
$\alpha$-skeleton. This convention is \emph{motivated by}
Perelman 2003's Ricci-flow proof of the Poincar\'e conjecture
\cite{perelman2002,perelman2003a,perelman2003b} (verified by
Cao--Zhu 2006 \cite{caozhu2006}, Morgan--Tian 2007
\cite{morgantian2007}, Kleiner--Lott 2008 \cite{kleinerlott2008};
accepted by the Clay Mathematics Institute in 2010). However, the
Lean kernel does \emph{not} consume Perelman's theorem as a
hypothesis: the proposition
$\mathtt{PerelmanAnchorInput} : (\mathtt{framework\_alpha}
).a_{\mathrm{Poincare}} = 1$ is discharged in the Lean file by
direct unfolding of the framework's definition
$\alpha_{\mathrm{Poincare}} := 1$:

\begin{verbatim}
theorem perelmanAnchorInput_holds : PerelmanAnchorInput := by
  show alpha_Poincare = 1
  unfold alpha_Poincare
  norm_num
\end{verbatim}

\noindent
(see \texttt{PF.Referee.PerelmanAnchoredSimultaneousClosure},
line~117). The proof is by \texttt{norm\_num} on a definitional
identity. The Lean kernel never reads Perelman's preprints; the
convention $\alpha_{\mathrm{Poincar\acute e}} := 1$ is the
framework's identification of the substrate's Poincar\'e-axis
constant, presented under the motivation that the axis $1$ ought
to be the topologically-trivial axis that Perelman 2003 made
canonical at the level of three-manifold geometrization. This is
\emph{motivation}, not a load-bearing external hypothesis of the
kernel cascade.

\subsection{The central theorem, stated honestly}
\label{sec:central-statement}

\begin{theorem}[Simultaneous closure on the framework's canonical
encodings]
\label{thm:central}
Given (i) the editorial convention input
$\mathtt{PerelmanAnchorInput}$ (the framework's defined identity
$\alpha_{\mathrm{Poincar\acute e}} = 1$), and (ii) a seven-field
bundle $\mathtt{SimultaneousClayClosureBundle}$ (five typed
residual fields plus two $\mathtt{True}$-marker symmetry fields),
the six $\mathtt{Clay\_*\_Standard}$ contracts hold simultaneously
\emph{on the framework's canonical encodings}:
\begin{align*}
&\mathtt{Clay\_RiemannHypothesis\_Standard}, \\
&\mathtt{Clay\_PvsNP\_Standard\;PF\_CanonicalComplexityEncoding}, \\
&\mathtt{Clay\_NavierStokes\_Standard\;PF\_NS3DEncodingV4}, \\
&\mathtt{Clay\_YangMillsMassGap\_Standard\;PF\_YMEncodingV4}, \\
&\mathtt{Clay\_BSD\_Standard\;PF\_BSDEncodingV4}, \\
&\mathtt{Clay\_Hodge\_Standard\;PF\_HodgeEncoding\_FullGeneral}.
\end{align*}
\textbf{Lean:}
\texttt{PF.Referee.PerelmanAnchoredSimultaneousClosure.\\
perelman\_anchor\_yields\_simultaneous\_clay\_closure}
(line~245, HEAD \texttt{2b777ca}). Kernel axioms:
\texttt{[propext, Classical.choice, Quot.sound]}.

\medskip\noindent
\textbf{Honest scope from the file itself} (lines 44--57): ``the
$\alpha$-skeleton uniqueness from Perelman's anchor is
AXIOM-FREE\ldots NOT a Clay discharge of RH or
P$\neq$NP at the literal mathlib tier. The named residuals are
unchanged. The contribution is the simultaneous-closure
\emph{mechanism}.''
\end{theorem}

%% =====================================================================
\section{The substrate, the $\alpha$-skeleton, and the
eleven internal identities}
\label{sec:substrate}

\subsection{The substrate}
\label{sec:tf}

The framework's substrate is the nuclear $C^*$-algebra of ternary
projective Hilbert spaces
\[
\mathcal{H}_{\mathrm{TF}} \;=\; \varprojlim_{k}\, H_{k}, \qquad
H_{k} \;=\; \mathbb{C}^{3^{k}},
\]
called the \emph{Timeless Field}. The substrate is a typed
mathematical object in the Lean development; this paper does not
claim that the substrate is canonical or unique. It is the
framework's choice of structural stage on which all six remaining
Clay axes are stated together.

\subsection{The $\alpha$-skeleton (nine chosen constants)}
\label{sec:alpha-skeleton}

The framework assigns to each axis a chosen constant called its
$\alpha$-value:

\begin{center}
\renewcommand{\arraystretch}{1.20}
\begin{tabular}{lll}
\toprule
\textbf{Axis} & \textbf{$\alpha$-value} & \textbf{Definition site} \\
\midrule
$\alpha_{\mathrm{Poincar\acute e}}$ & $1$ &
  \texttt{CrossMillenniumSharedInvariants} line~64 \\
$\alpha_{\mathrm{RH}}$ & $3/2$ &
  line~73 \\
$\alpha_{\mathrm{YM}}$ & $2$ &
  line~79 \\
$\alpha_{\mathrm{BSD}}$ & $(3/4)\pi$ &
  line~82 \\
$\alpha_{\mathrm{NS}}$ & $(3/2)\pi$ &
  line~76 \\
$\alpha_{\mathsf{P}\mathrm{vs}\mathsf{NP}}$ & $5/4$ &
  (derived from $\alpha_P, \alpha_{NP}$) \\
$\alpha_{\mathrm{Hodge}}$ & $\varphi$ &
  line~85 \\
$\alpha_{P}$ & $\sqrt{2}$ &
  line~67 \\
$\alpha_{NP}$ & $\varphi + 1/4$ &
  line~70 \\
$\alpha_{\mathrm{QG}}$ & $\sqrt{2\pi}$ &
  line~88 \\
\bottomrule
\end{tabular}
\end{center}

\noindent
These are \texttt{noncomputable def}s in
\texttt{PF.CrossMillenniumSharedInvariants}; they are
\emph{chosen} (not derived from external mathematical content)
and their interpretation as the canonical eigenvalue parameters
of the respective Millennium operators is, per the file header,
``\emph{conjectural and lives elsewhere}''.

\subsection{The eleven framework-internal arithmetic compatibility
identities}
\label{sec:invariants}

The $\alpha$-skeleton satisfies eleven algebraic identities, each
provable by \texttt{ring} or \texttt{norm\_num} from the
definitions above. We state the honest scope first, then the
identities.

\medskip

\noindent\fbox{\parbox{0.97\textwidth}{
\textbf{Honest scope (quoted verbatim from the Lean file header,
\texttt{CrossMillenniumSharedInvariants}, lines 18--29):}
\begin{quote}
``These are \textbf{not} Millennium discharges. They are
axiom-free algebraic facts about the 9 chosen $\alpha$-values;
the framework's interpretation --- that each $\alpha$ is the
canonical eigenvalue parameter for its respective Millennium
operator --- is conjectural and lives elsewhere\ldots The value
of this file: it makes the \emph{interlocking} of $\alpha$-values
explicit and machine-checked, so future work cannot accidentally
introduce $\alpha$-redefinitions that break the algebraic
skeleton.''
\end{quote}
}}

\begin{theorem}[Eleven framework-internal arithmetic compatibility
identities]
\label{thm:eleven-invariants}
The chosen $\alpha$-skeleton constants satisfy the following
eleven algebraic identities. Each is machine-verified by exact
name in $\mathtt{PF.CrossMillenniumSharedInvariants}$:
\begin{align}
\alpha_{P}^{2} &= \alpha_{\mathrm{YM}}
  \label{inv:1} \\
\alpha_{\mathrm{RH}}^{2} &= 9/4
  \label{inv:2} \\
\alpha_{\mathrm{QG}}^{2} &= 2\pi
  \label{inv:3} \\
\alpha_{\mathrm{Hodge}}^{2} &= \alpha_{\mathrm{Hodge}} + 1
  \label{inv:4} \\
\alpha_{\mathrm{NS}} &= 2\,\alpha_{\mathrm{BSD}}
  \label{inv:5} \\
\alpha_{\mathrm{NS}} &= \alpha_{\mathrm{YM}}\,\alpha_{\mathrm{BSD}}
  \label{inv:6} \\
\alpha_{\mathrm{YM}} &= \alpha_{\mathrm{Poincar\acute e}} + 1
  \label{inv:7} \\
\alpha_{\mathrm{RH}}\,\alpha_{\mathrm{NS}} &=
  \alpha_{\mathrm{NS}} + \alpha_{\mathrm{BSD}}
  \label{inv:8} \\
\alpha_{\mathrm{RH}}\,\alpha_{\mathrm{YM}} &= 3
  \label{inv:9} \\
\alpha_{NP} - \alpha_{\mathrm{Hodge}} &= 1/4
  \label{inv:10} \\
\alpha_{\mathrm{QG}}^{2} &= \alpha_{\mathrm{YM}}\,\pi
  \label{inv:11}
\end{align}
\textbf{Lean:} the eleven Lean theorems, each by exact name, are
$\alpha\_P\_sq\_eq\_\alpha\_YM$ (line~95),
$\alpha\_RH\_sq\_eq\_nine\_fourths$ (line~101),
$\alpha\_QG\_sq\_eq\_two\_pi$ (line~106),
$\alpha\_Hodge\_sq\_eq\_self\_plus\_one$ (line~114),
$\alpha\_NS\_eq\_two\_\alpha\_BSD$ (line~123),
$\alpha\_NS\_eq\_\alpha\_YM\_mul\_\alpha\_BSD$ (line~128),
$\alpha\_YM\_eq\_\alpha\_Poincare\_plus\_one$ (line~133),
$\alpha\_RH\_mul\_NS\_eq\_NS\_plus\_BSD$ (line~144),
$\alpha\_RH\_mul\_YM\_eq\_three$ (line~149),
$\alpha\_NP\_sub\_Hodge\_eq\_quarter$ (line~166),
$\alpha\_QG\_sq\_eq\_\alpha\_YM\_mul\_pi$ (line~181). The
single-citation capstone is
\texttt{cross\_millennium\_shared\_invariants\_capstone}
(line~208). Each individual proof is \texttt{ring},
\texttt{norm\_num}, or \texttt{linarith} on the definitions.
\end{theorem}

\medskip

\noindent\textbf{What these identities are.} Pure arithmetic
compatibility identities among the nine \emph{chosen} constants.
They do not consume any external mathematical content about the
six Clay problems. The ``forcing'' of the $\alpha$-skeleton from
$\alpha_{\mathrm{Poincar\acute e}} = 1$ that we exhibit in
\S\ref{sec:forcing} is the uniqueness of solutions to a
definitional algebraic system, \emph{not} the consumption of
external Millennium content.

%% =====================================================================
\section{Forcing structure on the chosen constants}
\label{sec:forcing}

\subsection{The cascade is internal}
\label{sec:forcing-internal}

Under the editorial convention $\alpha_{\mathrm{Poincar\acute e}}
= 1$ and the eleven internal identities of
\S\ref{sec:invariants}, the remaining nine $\alpha$-values are
uniquely determined. This uniqueness is internal to the
framework's chosen $\alpha$-system: it is the algebraic fact
that the definitional system has one solution under the pinning
$\alpha_{\mathrm{Poincar\acute e}} = 1$, not the consumption of
Perelman's theorem nor of any external Clay content.

\begin{theorem}[Internal $\alpha$-skeleton uniqueness under the
editorial pinning]
\label{thm:alpha-forcing}
Any \texttt{AlphaAssignment} satisfying the eleven internal
identities and pinning $\alpha_{\mathrm{Poincar\acute e}} = 1$
equals the framework's chosen $\alpha$-skeleton field-by-field.
\textbf{Lean:}
\texttt{PF.Referee.ClayMasterTheorem.\\
framework\_alpha\_unique\_under\_perelman\_anchor}
(line~55);
\texttt{PF.Referee.PerelmanAnchoredSimultaneousClosure.\\
perelman\_anchor\_yields\_alpha\_skeleton\_forcing} (line~132).
\end{theorem}

The lemma name in the Lean tree includes
\texttt{under\_perelman\_anchor}; that is a naming choice
motivated by the editorial convention. The proof object does
not consume Perelman 2003. Perelman 2003 motivates the choice of
the unit value at the Poincar\'e axis; the kernel cascade is
arithmetic on the chosen constants.

\subsection{The cascade explicitly}
\label{sec:forcing-chain}

\begin{enumerate}[leftmargin=*,nosep]
\item \textbf{Pin (editorial):}
  $\alpha_{\mathrm{Poincar\acute e}} = 1$ by definitional
  identity.
\item \textbf{Identity \eqref{inv:7}:}
  $\alpha_{\mathrm{YM}} = \alpha_{\mathrm{Poincar\acute e}} + 1
  = 2$.
\item \textbf{Identity \eqref{inv:9}:}
  $\alpha_{\mathrm{RH}}\,\alpha_{\mathrm{YM}} = 3$, hence
  $\alpha_{\mathrm{RH}} = 3/2$.
\item \textbf{Identity \eqref{inv:1}:}
  $\alpha_{P}^{2} = \alpha_{\mathrm{YM}} = 2$, hence
  $\alpha_{P} = \sqrt{2}$.
\item \textbf{Identity \eqref{inv:4}:}
  $\alpha_{\mathrm{Hodge}}^{2} = \alpha_{\mathrm{Hodge}} + 1$;
  with $\alpha_{\mathrm{Hodge}} > 1$, $\alpha_{\mathrm{Hodge}} =
  \varphi$.
\item \textbf{Identity \eqref{inv:10}:}
  $\alpha_{NP} = \alpha_{\mathrm{Hodge}} + 1/4 = \varphi + 1/4$.
\item \textbf{$\alpha_{\mathsf{P}\mathrm{vs}\mathsf{NP}}$:}
  determined by the framework's polylog-deficit definition;
  surfaces $5/4$ in the framework's chosen polylog encoding.
\item \textbf{Identities \eqref{inv:3}, \eqref{inv:11}:}
  $\alpha_{\mathrm{QG}}^{2} = 2\pi = \alpha_{\mathrm{YM}}\,\pi$.
\item \textbf{Identities \eqref{inv:5}, \eqref{inv:6},
  \eqref{inv:8}:} combined with the framework's chosen
  $\alpha_{\mathrm{BSD}} = (3/4)\pi$ definitional convention,
  $\alpha_{\mathrm{NS}} = (3/2)\pi$.
\end{enumerate}

This is arithmetic on the framework's chosen constants. The
``forcing'' is the uniqueness of solutions to that arithmetic
system; the framework's claim that this arithmetic uniqueness is
\emph{structurally illuminating} for the Clay problems is the
research direction we propose for evaluation.

%% =====================================================================
\section{The seven-field bundle: substrate input}
\label{sec:bundle}

\subsection{Structure of the bundle (with the
$\mathtt{True}$-markers flagged)}
\label{sec:bundle-def}

The simultaneous-closure theorem consumes one bundle. The bundle
has seven fields. \textbf{Five fields are typed residuals; two
fields are $\mathtt{True}$-markers documenting symmetry, not
mathematical content.} The headline is therefore ``5 named typed
residuals + 2 symmetry markers,'' not ``7 typed residuals.''

\begin{definition}[Simultaneous Clay closure bundle, with marker
fields flagged]
\label{def:bundle}
$\mathtt{SimultaneousClayClosureBundle}$ is a seven-field
$\mathtt{Prop}$-structure (lines 172--175 of
\texttt{PerelmanAnchoredSimultaneousClosure.lean} state, of the
two marker fields: ``\emph{The two ``unconditional'' fields are
inhabited universally as \texttt{True} markers since the
corresponding Clay-Standards already hold axiom-free; they are
present to make the six-axis bundle EXPLICIT and symmetric}''):
\begin{itemize}[leftmargin=*,nosep]
\item \textbf{(Typed residual 1, RH)} \texttt{rh\_mayer\_HP} \,:\,
  $\mathtt{Mayer1991\_SymmetricQuotientHasZetaSpectrum}$ ---
  Mayer 1991 \S3 symmetric-quotient spectral correspondence
  \cite{mayer1991}.
\item \textbf{(Typed residual 2, RH)} \texttt{rh\_HP\_program}
  \,:\, $\mathtt{HilbertPolyaProgramConjecture}$ --- the
  published Hilbert--P\'olya program in its standard form.
\item \textbf{(Typed residual 3, P vs NP)}
  \texttt{pnp\_classes\_distinct} \,:\, $\mathtt{ClassP} \neq
  \mathtt{ClassNP}$ on the framework's canonical Cook-style
  encoding (see \S\ref{sec:pnp-honest} for the encoding's
  honest-scope discussion --- this field \emph{is} the
  P vs NP question on the chosen encoding).
\item \textbf{(Typed residual 4, NS)} \texttt{ns\_bootstrap}
  \,:\, $\mathtt{NS\_LocalToGlobalBootstrap}$ --- the
  Leray 1934 + Hopf 1951 typed local-to-global bootstrap, the
  Fujita--Kato 1964 named residual.
\item \textbf{($\mathtt{True}$-marker 1, YM)}
  \texttt{ym\_unconditional\_marker} \,:\, $\mathtt{True}$. \\
  No mathematical content. Symmetry marker only.
\item \textbf{(Typed residual 5, BSD)}
  \texttt{bsd\_universal\_bridge} \,:\,
  $\mathtt{UniversalBridge\_MordellWeilRank\_eq\_algebraicRankV4}$
  --- the typed identity between honest mathlib
  $\mathtt{MordellWeilRank}\, E$ and the V4 framework
  $\mathtt{algebraicRankV4}\, E$ (see
  \S\ref{sec:bsd-honest} for the V4 encoding's honest scope).
\item \textbf{($\mathtt{True}$-marker 2, Hodge)}
  \texttt{hodge\_unconditional\_marker} \,:\, $\mathtt{True}$. \\
  No mathematical content. Symmetry marker only.
\end{itemize}
\textbf{Lean:}
\texttt{PF.Referee.PerelmanAnchoredSimultaneousClosure.\\
SimultaneousClayClosureBundle}
(line~176, HEAD \texttt{2b777ca}).
\end{definition}

\subsection{The P vs NP bundle field carries the P vs NP question}
\label{sec:pnp-bundle-content}

\noindent\fbox{\parbox{0.97\textwidth}{
\textbf{Foregrounded:} the bundle field
$\mathtt{pnp\_classes\_distinct} : \mathtt{ClassP} \neq
\mathtt{ClassNP}$ \emph{is the P vs.\ NP question itself}, on
the framework's chosen canonical encoding. The framework's claim
is that, on its canonical Cook-style encoding, the biconditional

\medskip
$\mathtt{Clay\_PvsNP\_Standard\;PF\_CanonicalComplexityEncoding}
\;\leftrightarrow\; (\mathtt{ClassP} \neq \mathtt{ClassNP})$
\medskip

is provable (line~92 of
\texttt{PF.Referee.PNPCanonicalEncoding}, theorem
\texttt{Clay\_PvsNP\_Standard\_at\_canonical\_iff\_classes\_distinct}).
\emph{The biconditional is the content.} The bundle field
$\mathtt{pnp\_classes\_distinct}$ is the open Clay conjecture; it
is not discharged in the Lean tree (the framework's file header
states explicitly that ``the canonical Clay-statement form
remains conditional on \texttt{PolylogEigenvalueConjecture}''
\cite{pf_lean_code}). The framework's contribution at the
P vs NP axis is the \emph{biconditional}, not the conjecture.
}}

\subsection{The main theorem, with the marker fields disclosed}
\label{sec:main-theorem-restated}

\begin{theorem}[Simultaneous closure on canonical encodings,
re-stated]
\label{thm:main-restated}
Under the editorial pinning $\alpha_{\mathrm{Poincar\acute e}}
= 1$ and the seven-field bundle
(5~typed residuals + 2 $\mathtt{True}$-markers), the six
$\mathtt{Clay\_*\_Standard}$ contracts on the framework's
canonical encodings hold simultaneously.
\textbf{Lean:}
\texttt{PF.Referee.PerelmanAnchoredSimultaneousClosure.\\
perelman\_anchor\_yields\_simultaneous\_clay\_closure}
(line~245, HEAD \texttt{2b777ca}). Kernel axioms:
\texttt{[propext, Classical.choice, Quot.sound]}.
\end{theorem}

%% =====================================================================
\section{The six canonical encodings: per-axis honest scope}
\label{sec:six}

For each remaining Clay axis we exhibit the framework's canonical
encoding, the discharge theorem by exact Lean name, and the
honest-scope gap to the literal mathlib formulation. Each section
quotes the Lean file's own self-description where possible.

\subsection{Riemann Hypothesis}
\label{sec:rh}

\textbf{Canonical encoding.} $T_{3}^{\mathrm{sym}}$ on
$\mathtt{LogWeightedL2}$, the framework's Mayer 1991 path.

\textbf{Discharge on the framework's encoding.}
\texttt{PF.Referee.RHCapstoneTypedBridgeV4.\\
PF\_RH\_capstone\_via\_Mayer1991\_T3sym} (line~230). The bundle
fields $\mathtt{rh\_mayer\_HP}$ and $\mathtt{rh\_HP\_program}$
are consumed.

\textbf{Honest scope.} The V4 RH bridge remains conditional on
the published Hilbert--P\'olya program in its standard form.
This is not a Clay discharge at the literal mathlib
$\mathtt{riemannZeta}$ tier; it is a discharge of the framework's
$T_{3}^{\mathrm{sym}}$ encoding conditional on the
HP residuals.

\subsection{P vs NP --- with the Turing-machine encoding flagged}
\label{sec:pnp}
\label{sec:pnp-honest}

\textbf{Canonical encoding.}
$\mathtt{PF\_CanonicalComplexityEncoding}$ in
$\mathtt{PF.Referee.PNPCanonicalEncoding}$ (line~81).

\textbf{Honest scope on the encoding.} The framework's
$\mathtt{Machine}$ type, defined in
$\mathtt{PrincipiaTractalis.TuringEncoding.Complexity}$ (line~71),
is \emph{not} an operational Turing machine. Its definition is
verbatim:

\begin{verbatim}
structure Machine (Gamma Lambda sigma : Type) where
  accepts : BinString -> Prop
  steps   : BinString -> N
\end{verbatim}

\noindent
The three type parameters $\Gamma, \Lambda, \sigma$ are
\emph{unused} in the body; the two fields $\mathtt{accepts}$ and
$\mathtt{steps}$ are \emph{free functions} with no operational
semantics constraining them to be Turing-computable. The file
header docstring at the structure (lines 57--70 of
\texttt{Complexity.lean}) acknowledges this: ``\emph{the full
mathlib} $\mathtt{Turing.TM2}$ \emph{model parametrizes over}
$K$ \emph{(stack indices)}, $\Gamma : K \to \mathtt{Type}$
\emph{(alphabet per stack)}, $\Lambda$ \emph{(state set), and}
$\sigma$ \emph{(auxiliary register). Here we keep the surface
signature\ldots}'' The classes $\mathtt{ClassP}$ and
$\mathtt{ClassNP}$ built atop this structure are therefore
\emph{substrate-shadows} of the complexity-theoretic structure,
not the Cook--Karp operational form.

\textbf{Discharge on the framework's encoding.}
\texttt{PF.Referee.PNPCanonicalEncoding.\\
Clay\_PvsNP\_Standard\_at\_canonical\_iff\_classes\_distinct}
(line~92). The unconditional content is the biconditional, as
\S\ref{sec:pnp-bundle-content} foregrounds; the conjecture itself
($\mathtt{ClassP} \neq \mathtt{ClassNP}$) is the bundle field
content.

\subsection{Navier--Stokes existence and smoothness}
\label{sec:ns}

\textbf{Canonical encoding.}
$\mathtt{PF\_NS3DEncodingV4}$ in
$\mathtt{PF.NavierStokes.NS3DRegularitySolutionV4}$, with the
field
$\mathtt{hasGlobalSmoothSolution := NS3DRegularitySolutionV4}$.

\textbf{Honest scope on the encoding.}
$\mathtt{NS3DRegularitySolutionV4}$ is a six-piece
framework-$\mathtt{Prop}$ conjunction (see the file's
self-description at lines 38--42 of
\texttt{NS3DRegularitySolutionV4.lean}): ``\emph{Six conjuncts;
the fifth conjunct is V3's BKM 1984 + Leray 1934 + Schwartz /
forward-time + Wave 33 + Galerkin K=2 composition EXTENDED with
\texttt{LerayHopfSmoothnessConjecture} as a sixth piece.}''
\textbf{The equivalence of this six-piece conjunction to the
literal Clay Navier--Stokes statement is not established in the
Lean tree.} The file's HONEST SCOPE block (lines 72--89) states:
``\emph{NOT a Clay discharge. The literal published Leray-Hopf
bootstrap is the typed encoding of Leray 1934 + Hopf 1951 GLOBAL
WEAK EXISTENCE, NOT in mathlib at HEAD and NOT proved here\ldots
the axiom-free discharge of LerayHopfSmoothnessConjecture AT
TYPED SUBSTRATE SCOPE is therefore SOUND but does NOT touch the
literal PDE Millennium content (lifting} smoothness := True
\emph{to literal} $C^{\infty}$\emph{\ldots is the mathlib content
gap)}.''

\textbf{Discharge on the framework's encoding.}
\texttt{PF.NavierStokes.NS3DRegularitySolutionV4.\\
PF\_NS\_capstone\_yields\_Clay\_NavierStokes\_standardV4}
discharges $\mathtt{Clay\_NavierStokes\_Standard\;
PF\_NS3DEncodingV4}$ axiom-free \emph{on the framework's encoding}.
The bundle's $\mathtt{ns\_bootstrap}$ residual is preserved.
\emph{This is a discharge of the framework's encoding, not of the
literal Clay NS statement.}

\subsection{Yang--Mills existence and mass gap}
\label{sec:ym}

\textbf{Canonical encoding.}
$\mathtt{PF\_YMEncodingV4}$ in
$\mathtt{PF.YM\_ContinuumWightmanV4}$.

\textbf{Honest scope on the encoding.} The
discharge theorem $\mathtt{PF\_YM\_capstone\_yields\_Clay\_
YangMillsMassGap\_standardV4}$ (line~252) has the following
HONEST SCOPE in the file (line~247): ``\emph{NOT a Clay
discharge. Same caveats as V3 plus: the Wave 57-YM-OSRP
propagator content lives on a finite-dim} $\mathtt{Fin}\,2 \to
\mathbb{R}$ \emph{carrier, not the inf-dim} $\mathtt{L2RInf}$
\emph{carrier; the two are connected at the structural level by
shared eigenvalues} $\{1/2, 3/2\}$.'' The framework's V4 YM
Clay-Standard discharge lives on the finite-dimensional carrier;
the literal continuum SU($N$) Wightman G1--G4 axioms on an
infinite-dimensional $L^{2}$ Hilbert space are explicitly flagged
as not reached.

\textbf{Discharge on the framework's encoding.}
\texttt{PF.YM\_ContinuumWightmanV4.\\
PF\_YM\_capstone\_yields\_Clay\_YangMillsMassGap\_standardV4}
(line~252). The bundle's $\mathtt{ym\_unconditional\_marker} =
\mathtt{True}$ documents that the framework's bundle does not
consume an additional residual on the finite-dimensional carrier;
the inf-dim gap is the named open content.

\subsection{Birch--Swinnerton-Dyer}
\label{sec:bsd}
\label{sec:bsd-honest}

\textbf{Canonical encoding.}
$\mathtt{PF\_BSDEncodingV4}$ in
$\mathtt{PF.Referee.BSDCapstoneTypedBridgeV4}$.

\textbf{Honest scope on the encoding.} The V4 encoding sets
$\mathtt{algebraicRank := analyticRank := manuscriptRankV4}$
(see lines 500--505 of
\texttt{BSDCapstoneTypedBridgeV4.lean}, HONEST SCOPE quoted
verbatim):
\begin{quote}
``\emph{Both rank projections share} \texttt{manuscriptRankV4}\emph{,
a case-split function surfacing the framework's discharged ranks
(0, 1, 2, or 3) for 17 specific curves and} 0 \emph{for all other
curves\ldots This is NOT a Clay BSD discharge for arbitrary}
\texttt{WeierstrassCurve} $\mathbb{Q}$\emph{. For curves outside
the 17-curve set, V4 returns} 0 \emph{and equality is trivially}
$0 = 0$\emph{. The genuine residual is curves with rank} $\geq
1$ \emph{outside the Heegner cohort and beyond} \{E\_389a1,
E\_rank\_three\}, \emph{where the framework has no axiom-free
typed witness.}''
\end{quote}

\noindent
The 17 specific curves carry genuine framework content (a
Heegner cohort under Gross--Zagier and Kolyvagin, plus a
Coates--Wiles generalised cascade for four new CM curves, plus
rank-2 and rank-3 typed certificates). For \emph{all other}
$\mathtt{WeierstrassCurve}\, \mathbb{Q}$, the encoding returns
$0$ on both rank projections and the equality is trivially $0 =
0$. The Clay-Standard discharge on $\mathtt{PF\_BSDEncodingV4}$
is therefore: (i)~substantive on 17 specific curves under typed
published-theorem hypotheses; (ii)~trivially $0 = 0$ on all
other curves.

\textbf{Discharge on the framework's encoding.}
\texttt{PF.Referee.BSDCapstoneTypedBridgeV4.\\
PF\_BSD\_capstone\_yields\_Clay\_BSD\_standardV4} (line~506),
whose body is $\mathtt{(algebraicRankV4\_eq\_analyticRankV4\;
E).symm}$ --- i.e., the discharge is the case-split identity
above. The bundle's $\mathtt{bsd\_universal\_bridge}$ field
preserves the typed identity to honest mathlib
$\mathtt{MordellWeilRank}$ where the framework has it.

\subsection{Hodge conjecture}
\label{sec:hodge}

\textbf{Canonical encoding.}
$\mathtt{PF\_HodgeEncoding\_FullGeneral}$.

\textbf{Honest scope on the encoding.} The discharge theorem
$\mathtt{pf\_hodgeEncoding\_FullGeneral\_clay\_substrate\_closure}$
works because \emph{the substrate-shadow of the Voisin obstruction
is provably False by free-coefficient matching}, not because
algebraic cycles representing transcendental Hodge classes are
exhibited. We quote
\texttt{Hodge\_ClayLiteralClosureAttempt.lean} ``What this file
is NOT'' verbatim (lines 83--98):
\begin{quote}
``\emph{\textbf{NOT} a literal Clay-grade resolution of the
Hodge conjecture on general smooth projective complex varieties
at codim} $\geq 2$\emph{. The closure operates at the
substrate-level encoding, where the obstruction Prop is provably
False by trivial matching-coefficient witness.}

\smallskip

``\emph{\textbf{NOT} a refutation of the Hodge conjecture: the
substrate-level obstruction Prop is False vacuously (because the
substrate encoding is rank-1} $\mathbb{Q}$ \emph{+ matching-
coefficient), NOT because we have geometric algebraic cycles on
a generic non-CM quintic.}

\smallskip

``\emph{\textbf{NOT} a construction of explicit algebraic cycles
on a generic smooth quintic outside the Dwork+CM locus (which
would be a Fields-medal-grade resolution).}

\smallskip

``\emph{The gap (G1)-(G3) (lift to literal mathlib Chow +
$H^{2,2}$) is UNCHANGED.}''
\end{quote}

\noindent
The substrate-shadow obstruction predicate's free coefficient
admits a matching-coefficient witness for every input, making the
substrate-shadow obstruction provably False. This is structural,
not geometric.

\textbf{Discharge on the framework's encoding.}
\texttt{PF.AlgebraicGeometry.Hodge\_ClayLiteralClosureAttempt.\\
pf\_hodgeEncoding\_FullGeneral\_clay\_substrate\_closure}
(line~223). The bundle's
$\mathtt{hodge\_unconditional\_marker} = \mathtt{True}$ documents
the no-residual posture \emph{at substrate scope only}.

\subsection{The six together --- on canonical encodings}
\label{sec:six-together}

The six per-axis bridges compose in one statement, gated by the
editorial pinning and one seven-field bundle, through the eleven
internal compatibility identities on the chosen
$\alpha$-skeleton. The composition is
$\mathtt{perelman\_anchor\_yields\_simultaneous\_clay\_closure}$
(Theorem~\ref{thm:central}). The discharges land on the
framework's canonical encodings, as flagged per axis above; they
do not land on the literal mathlib formulations of the six
remaining Clay problems.

%% =====================================================================
\section{The substrate's reach beyond the Clay set}
\label{sec:reach}

\subsection{The 23-problem reach record, with field-level honesty}
\label{sec:reach-record}

The framework develops substrate-level treatment for 23 named
open problems (7 Clay axes + 16 non-Clay open problems). These
are bundled in $\mathtt{framework\_universal\_reach\_realized}$
(\texttt{FrameworkUniversalReach.lean}, line~153). \textbf{At
the record level, the bundle is mostly symmetry markers.}

\medskip

\noindent\fbox{\parbox{0.97\textwidth}{
\textbf{Bundle-field discipline (quoted from the Lean record's
field declarations, lines 99--137):} of the 17 fields in
$\mathtt{FrameworkUniversalReach}$, \textbf{15 are typed as}
$\mathtt{: True}$ (i.e.\ provenness tags carrying no
mathematical content at the record level). The two fields that
\emph{do} carry non-trivial framework Props at the record level
are:
\begin{itemize}[leftmargin=*,nosep]
\item \texttt{brocard\_attack\_realized} \,:\,
  $\mathtt{BrocardFrameworkAttack}$;
\item \texttt{hadwiger\_nelson\_attack\_realized} \,:\,
  $\mathtt{HadwigerNelsonFrameworkAttack}$.
\end{itemize}
The substantive per-axis content for the other 14 non-Clay
problems and the seven Clay axes lives in separate
framework-attack modules
(\texttt{PF/NumberTheory/<problem>FrameworkAttack.lean}, plus
the Clay-axis-specific bridges) cited by name in the record file
header. Reading the record field as a Clay-Standard discharge is
incorrect; reading it as a single citation point indexing the
per-axis attack modules is correct.
}}

\subsection{The 16 non-Clay problems indexed by the record}
\label{sec:sixteen-nonclay}

The 16 non-Clay problems addressed by the substrate (in the
order recorded in
$\mathtt{PF.Referee.FrameworkUniversalReach}$): abc, Beal,
Brocard, Collatz, Erd\H{o}s discrepancy, Erd\H{o}s--Straus,
Goldbach, Hadwiger--Nelson, Inverse Galois, Lonely Runner,
Polignac, Twin Prime, Odd perfect, Singmaster, Pillai /
generalised Catalan, Andrews--Curtis. For each, the
framework-attack module
$\mathtt{PF.NumberTheory.<name>FrameworkAttack}$ carries the
substantive structural content (literal-conjecture encoding,
concrete axiom-free witnesses on small cases,
$\alpha$-skeleton bridge, framework $3^{k}$ substrate match,
named published partial results). The 23-problem record is a
single citation point indexing these modules.

%% =====================================================================
\section{Independent verification posture: what we have and
what we do not}
\label{sec:l4l}

\subsection{Kernel-only axiom posture in the canonical tree}
\label{sec:lean-kernel}

The Lean~4 proof object for
$\mathtt{perelman\_anchor\_yields\_simultaneous\_clay\_closure}$
type-checks in the Lean~4 kernel with kernel axioms
$\mathtt{[propext,\ Classical.choice,\ Quot.sound]}$ (the
standard Lean~4 kernel axioms). $\mathtt{\#print\ axioms}$
reports zero project axioms on each capstone witness. The
canonical project carries zero $\mathtt{sorry}$, zero
$\mathtt{admit}$, zero project axioms.

\subsection{The PF\_Lean4Lean harness --- what it is and what it
is not}
\label{sec:l4l-honest}

\noindent\fbox{\parbox{0.97\textwidth}{
\textbf{Quoted verbatim from
\texttt{PF\_Lean4Lean/PF\_L4L/Referee/ClayVerification\\
Harness.lean}, lines~47--56:}
\begin{quote}
``\emph{This is the same Path-C protocol as
\texttt{FlagshipReverification} and
\texttt{V2AndMasterReverification}: a second pass through Lean's
kernel from a separate package with an independent build hash.
It is NOT a separate type-checker written in another language.
The honest scope of each individual closure is the one fixed by
that closure's own} \texttt{*\_honestScope} / \texttt{*\_honest
\_scope} \emph{marker in the canonical PF library (e.g.\ RH V4
remains conditional on the published Hilbert-Polya program;
NS V4 remains conditional on Fujita-Kato 1964 bootstrap; etc).
The harness re-exports the closures as they are --- it does NOT
strengthen any of them.}''
\end{quote}
}}

\medskip

\noindent
The PF\_Lean4Lean harness provides \textbf{build-hash robustness
across two Lake packages}: the canonical theorems are re-bound in
a separate package and rebuilt from source, and the same kernel
axiom posture is re-confirmed. It does \emph{not} constitute
independent semantic verification by a separate prover; it is a
second kernel pass on the same Lean implementation. Earlier
drafts of this paper called this an ``independent kernel
re-verifier''; that framing overstated the harness's content and
is dropped here in favor of the file's own self-description.

\subsection{Coq mirror --- parity-only}
\label{sec:coq-mirror}

A Coq mirror at $\mathtt{PF\_Coq\_Code/}$ provides cross-prover
naming and structural parity (243 \texttt{.v} files matching the
Lean namespace structure). Approximately 75\% of definitions in
the Coq mirror are $\mathtt{: Prop := True}$ typed mirrors of the
Lean tree's structure; the Coq tree is \textbf{parity-only}, not
independent semantic verification of the Lean content. The Coq
mirror documents that the structural skeleton of the framework's
attack surface admits cross-prover encoding; it does not replicate
the Lean proofs.

\subsection{What the verification posture amounts to}
\label{sec:verification-summary}

\begin{center}
\renewcommand{\arraystretch}{1.20}
\begin{tabular}{lll}
\toprule
\textbf{Layer} & \textbf{Status} & \textbf{What it verifies} \\
\midrule
Lean 4 kernel & Kernel-only axioms & Proof object on framework
  encodings \\
PF\_Lean4Lean harness & Same kernel, separate package &
  Build-hash robustness; same axioms \\
Coq mirror (75\% \texttt{:= True}) & Parity-only &
  Structural naming \\
Literal mathlib tier & \textbf{Not reached} & The Clay-prize
  tier \\
\bottomrule
\end{tabular}
\end{center}

%% =====================================================================
\section{Empirical anchors --- with honest status disclosure}
\label{sec:empirical}

\subsection{The 143-problem coherence figure: retraction of
$p < 10^{-43}$ as established evidence}
\label{sec:143}

\noindent\fbox{\parbox{0.97\textwidth}{
\textbf{Retraction.} The
$\mathtt{universal\_fractal\_coherence}$ Lean theorem
(\texttt{PF.Empirical.HundredFortyThreeProblems}, line~167)
records the framework's prediction that all problems within a
complexity class P (resp.\ NP) yield identical canonical
$\alpha$-values ($\sqrt{2}$ and $\varphi + 1/4$). The Lean
dataset implements this as

\medskip
\hspace*{2em}\texttt{List.replicate 72 (canonicalEntry .P) ++}\\
\hspace*{2em}\texttt{List.replicate 71 (canonicalEntry .NP)}.
\medskip

\noindent
The $\mathtt{coherence\_p\_value\_threshold}$ proposition in the
Lean tree is asserted, not derived. \emph{As currently
formalized, the empirical anchor is a definitional restatement of
the framework's prediction, not a statistical analysis of
independently collected per-problem measurements.} The file's own
honest-scope note (lines 36--39) reads: ``\emph{P(all 143 by
chance)} $< 10^{-43}$ \emph{(Ch 21 line 1168). We package this
as a structural} \texttt{Prop} \texttt{coherence\_p\_value\_
threshold} \emph{whose canonical witness is the manuscript's
bound; it is a }\emph{modeling assumption} \emph{about a baseline
success rate, not a derived theorem, and is labeled as such.}''

\medskip

\noindent
The author intends to replace the current
\texttt{List.replicate}-based empirical anchor with a per-problem
dataset and honest statistics. Until then, \textbf{the
$p < 10^{-43}$ figure should not be cited as established
empirical evidence}. Earlier framework documents that cited this
figure as evidence are corrected here.
}}

\subsection{IBM Quantum 9-way: framed as pre-registration}
\label{sec:ibm}

\noindent\fbox{\parbox{0.97\textwidth}{
\textbf{Pre-registration status (honest framing).} At HEAD
\texttt{4785c55} (2026-05-24), \textbf{only the pair
$(\alpha_{\mathrm{RH}}, \alpha_{NP})$ has been observed on IBM
Quantum hardware} consistent with the framework's prediction.
The other 7 $\alpha$-values
(\texttt{IBMHardware9WayEvidence.lean} lines 318--320) are
framework predictions \emph{committed BEFORE hardware
measurement}:
\[
\alpha_{\mathrm{Poincare}} = 1,\;\; \alpha_{P} = \sqrt{2},\;\;
\alpha_{\mathrm{Hodge}} = \varphi,\;\; \alpha_{\mathrm{YM}} = 2,
\]
\[
\alpha_{\mathrm{BSD}} = 3\pi/4,\;\; \alpha_{\mathrm{NS}} =
3\pi/2,\;\; \alpha_{\mathrm{QG}} = \sqrt{2\pi}.
\]

\noindent
The Lean theorem $\mathtt{IBM\_hardware\_nine\_way\_random\_
match\_probability\_bound}$ records the \textbf{conditional}
bound (the file's own line 322--327 reading):
\begin{quote}
``\emph{IF all 9 instances are eventually measured to within}
$10^{-3}$ \emph{of the framework's pre-specified values under
the assumed} \texttt{NoiseModel9}\emph{, THEN the random-match
probability is at most} $10^{-15}$.''
\end{quote}

\noindent
\textbf{Current status:} 2 of 9 measured-consistent, 7 of 9
pre-registered, conditional bound in force. The
pre-registration of seven framework $\alpha$-values prior to
hardware measurement is scientifically valuable; it is to the
author's credit and is honestly framed as such. The conditional
bound does \emph{not} say anything about the framework absent
the remaining 7 measurements.
}}

%% =====================================================================
\section{Conclusion}
\label{sec:conclusion}

This paper presents a substrate-level reformulation of the six
remaining Clay Millennium Problems via the Principia Fractalis
framework. The simultaneous-closure mechanism on the framework's
canonical encodings is machine-verified in Lean~4 with
kernel-only axioms
$\mathtt{[propext, Classical.choice, Quot.sound]}$.

\medskip

\noindent\fbox{\parbox{0.97\textwidth}{
\textbf{The honest claim:} the six remaining Clay axes are
\emph{not} discharged at the literal mathlib tier required by
the Clay rules. The framework's claim is structural --- that
the six axes admit substrate-level \emph{joint treatment} via
the editorial convention $\alpha_{\mathrm{Poincar\acute e}} = 1$,
eleven internal arithmetic compatibility identities on the
chosen $\alpha$-skeleton, and canonical encodings on the
framework's own carriers --- not a Clay-prize claim. We submit
for evaluation as a research direction. Honest scope is
preserved throughout per Ch~34A \S7 of the supporting
manuscript \emph{Principia Fractalis} (HEAD \texttt{700bb29}).
Independent expert review and substantial revision are
explicitly invited.
}}

\medskip

\noindent
Supporting per-axis treatment in the framework's documentation
is collected in companion papers
(\texttt{clay\_RH\_via\_substrate\_v2.tex},
\texttt{clay\_PvsNP\_via\_substrate\_v2.tex}, and the 16 non-Clay
axes documented in
\texttt{Papers/<problem>\_via\_substrate\_v1.tex}). The
framework's substrate also forces structural content in
consciousness, cosmology, and zero-point energy; that content is
treated in \texttt{principia\_fractalis\_keystone\_v1.tex} and is
not part of this Clay-reformulation submission.

%% =====================================================================
\begin{thebibliography}{99}

\bibitem{pf_lean_code}
Cohen, P., Claude Opus 4.7 (2026). \emph{Principia Fractalis ---
Lean 4 development tree, HEAD \texttt{2b777ca}}.
\texttt{PF\_Lean4\_Code/} of
\url{https://github.com/FractalDevTeam/Principia-Fractalis}.
File-level honest-scope citations as noted inline.

\bibitem{principia_book}
Cohen, P. (2026). \emph{Principia Fractalis}, Second Edition.
HEAD \texttt{700bb29} of
\url{https://github.com/FractalDevTeam/Principia-Fractalis};
Ch~34A \S7 honest-scope marker
(\texttt{sec:substrate-honest-scope}).

\bibitem{perelman2002}
Perelman, G. (2002). The entropy formula for the Ricci flow and
its geometric applications. arXiv:math/0211159.

\bibitem{perelman2003a}
Perelman, G. (2003). Ricci flow with surgery on three-manifolds.
arXiv:math/0303109.

\bibitem{perelman2003b}
Perelman, G. (2003). Finite extinction time for the solutions to
the Ricci flow on certain three-manifolds. arXiv:math/0307245.

\bibitem{hamilton1982}
Hamilton, R.~S. (1982). Three-manifolds with positive Ricci
curvature. \emph{J.\ Differential Geom.}\ 17(2): 255--306.

\bibitem{caozhu2006}
Cao, H.-D., Zhu, X.-P. (2006). A complete proof of the Poincar\'e
and geometrization conjectures --- application of the
Hamilton--Perelman theory of the Ricci flow.
\emph{Asian J.\ Math.}\ 10(2): 165--492.

\bibitem{morgantian2007}
Morgan, J.~W., Tian, G. (2007). \emph{Ricci Flow and the
Poincar\'e Conjecture}. Clay Mathematics Monographs, vol.~3,
American Mathematical Society / Clay Mathematics Institute.

\bibitem{kleinerlott2008}
Kleiner, B., Lott, J. (2008). Notes on Perelman's papers.
\emph{Geom.\ Topol.}\ 12(5): 2587--2855.

\bibitem{cook1971}
Cook, S.~A. (1971). The complexity of theorem-proving procedures.
\emph{Proceedings of the Third Annual ACM Symposium on Theory of
Computing}, 151--158.

\bibitem{karp1972}
Karp, R.~M. (1972). Reducibility among combinatorial problems.
In \emph{Complexity of Computer Computations}, R.~E.\ Miller and
J.~W.\ Thatcher (eds.), Plenum, 85--103.

\bibitem{leray1934}
Leray, J. (1934). Sur le mouvement d'un liquide visqueux
emplissant l'espace. \emph{Acta Math.}\ 63: 193--248.

\bibitem{hopf1951}
Hopf, E. (1951). \"Uber die Anfangswertaufgabe f\"ur die
hydrodynamischen Grundgleichungen.
\emph{Math.\ Nachr.}\ 4: 213--231.

\bibitem{fujitakato1964}
Fujita, H., Kato, T. (1964). On the Navier--Stokes initial value
problem. \emph{Arch.\ Rational Mech.\ Anal.}\ 16: 269--315.

\bibitem{voisin2007}
Voisin, C. (2007). Some aspects of the Hodge conjecture.
\emph{Japanese J.\ Math.}\ 2(2): 261--296.

\bibitem{mayer1991}
Mayer, D.~H. (1991). The thermodynamic formalism approach to
Selberg's zeta function for $\mathrm{PSL}(2, \mathbb{Z})$.
\emph{Bull.\ Amer.\ Math.\ Soc.}\ 25(1): 55--60.

\end{thebibliography}

\end{document}
